\documentclass{beamer}
\usetheme{Madrid}

\usepackage[utf8]{inputenc} % solo si usas pdflatex
\usepackage[T1]{fontenc}    % solo si usas pdflatex

\usepackage{amsmath,amssymb}
\usepackage{microtype}
\usepackage{cancel}         % si lo usas
\usepackage{fontawesome}    % si lo usas

\usepackage{hyperref}

\usepackage{graphicx}
\graphicspath{{images/}}

\hypersetup{pdfborder={0 0 0}}

\title{Tipos de BBDDs}
\author{Rodrigo Blanco Betes. Alfa Formación}
\date{Febrero del 2026}

\begin{document}

% --------- Portada ---------
\begin{frame}
  \titlepage
\end{frame}

% --------- Índice ---------
\begin{frame}{Contenido}
  \tableofcontents
\end{frame}

% =========================================================
\section{SQLite}

\begin{frame}{SQLite: ¿qué es?}
\begin{itemize}
  \item Base de datos relacional \textbf{embebida} (archivo único).
  \item Cero servidor: ideal para \textbf{apps}, \textbf{IoT}, \textbf{prototipos}, \textbf{tests}.
  \item Transaccional (ACID), muy portable y con huella mínima.
\end{itemize}
\end{frame}

\begin{frame}{SQLite: casos de uso típicos}
\begin{itemize}
  \item Aplicaciones móviles/desktop (caché local, almacenamiento offline).
  \item Sistemas embebidos/IoT (telemetría, configuración).
  \item Prototipado rápido y pruebas unitarias/integración.
  \item Data analysis ``ligero'' (ETL pequeño, staging local).
\end{itemize}
\end{frame}

\begin{frame}{SQLite: programas y comandos clave}
\textbf{Programas/herramientas:}
\begin{itemize}
  \item \texttt{sqlite3} (CLI principal)
  \item Navegadores GUI (p.ej. DB Browser for SQLite, DBeaver)
\end{itemize}
\end{frame}

\begin{frame}[fragile]{SQLite: programas y comandos clave}
\textbf{Programas/herramientas:}
\begin{itemize}
  \item \texttt{sqlite3} (CLI principal)
  \item Navegadores GUI (p.ej. DB Browser for SQLite, DBeaver)
\end{itemize}

\textbf{Comandos útiles (CLI \texttt{sqlite3}):}
\begin{verbatim}
sqlite3 mi_bd.sqlite
.help
.tables
.schema tabla
.headers on
.mode column
.timer on
.quit
\end{verbatim}
\end{frame}

\begin{frame}[fragile]{SQLite: SQL esencial (ejemplos)}
\begin{verbatim}
-- Crear tabla
CREATE TABLE usuarios (
  id INTEGER PRIMARY KEY,
  nombre TEXT NOT NULL,
  email TEXT UNIQUE
);

-- Insertar y consultar
INSERT INTO usuarios(nombre, email)
VALUES ('Ana','ana@ejemplo.com');

SELECT * FROM usuarios
WHERE email LIKE '%@ejemplo.com';

-- Indices y transacciones
CREATE INDEX idx_usuarios_email ON usuarios(email);

BEGIN;
UPDATE usuarios
SET nombre='Ana Gomez'
WHERE id=1;
COMMIT;
\end{verbatim}
\end{frame}

% =========================================================
\section{Neo4j}

\begin{frame}{Neo4j: ¿qué es?}
\begin{itemize}
  \item Base de datos de \textbf{grafos} (property graph).
  \item Optimizada para relaciones profundas y consultas de caminos.
  \item Lenguaje principal: \textbf{Cypher}.
\end{itemize}
\end{frame}

\begin{frame}{Neo4j: casos de uso típicos}
\begin{itemize}
  \item \textbf{Recomendadores}: ``usuarios que vieron X también vieron Y''.
  \item \textbf{Fraude} y detección de anomalías (redes de transacciones).
  \item \textbf{Knowledge graphs} (ontologías, entidades y relaciones).
  \item \textbf{Redes}: telecom, rutas, dependencias de microservicios.
  \item \textbf{Impact analysis}: propagación de cambios (dependencias).
\end{itemize}
\end{frame}

\begin{frame}[fragile]{Neo4j: programas y comandos clave}
\textbf{Programas/herramientas:}
\begin{itemize}
  \item \texttt{neo4j} (gestión del servidor)
  \item \texttt{cypher-shell} (CLI para Cypher)
  \item Neo4j Browser (UI web) y Neo4j Desktop (entorno local)
\end{itemize}

\textbf{Comandos típicos:}
\begin{verbatim}
neo4j start
neo4j status
neo4j stop

cypher-shell -u neo4j -p <password>
\end{verbatim}
\end{frame}

\begin{frame}[fragile]{Neo4j: Cypher esencial (ejemplos)}
\begin{verbatim}
-- Crear nodos y relaciones
CREATE (a:Persona {nombre:'Ana'})
CREATE (b:Persona {nombre:'Luis'})
CREATE (a)-[:AMIGO_DE]->(b);

-- Consultar patrones
MATCH (p:Persona)-[:AMIGO_DE]->(q:Persona)
RETURN p.nombre, q.nombre;

-- Busqueda de caminos
MATCH p = (a:Persona {nombre:'Ana'})-[:AMIGO_DE*1..3]->(x)
RETURN p;
\end{verbatim}
\end{frame}

\begin{frame}{Neo4j: buenas prácticas rápidas}
\begin{itemize}
  \item Modela \textbf{relaciones como primera clase} (evita ``joins manuales'').
  \item Crea \textbf{constraints} e índices para nodos clave (p.ej. \texttt{Persona(nombre)}).
  \item Mantén relaciones con semántica clara: \texttt{COMPRO}, \texttt{SIGUE\_A}, \texttt{DEPENDE\_DE}.
\end{itemize}
\end{frame}

% =========================================================
\section{MongoDB}

\begin{frame}{MongoDB: ¿qué es?}
\begin{itemize}
  \item Base de datos \textbf{documental} (BSON/JSON).
  \item Flexible: esquema evolutivo, ideal para datos semi-estructurados.
  \item Buen soporte para escalado horizontal (sharding) y replicación.
\end{itemize}
\end{frame}

\begin{frame}{MongoDB: casos de uso típicos}
\begin{itemize}
  \item Catálogos (e-commerce), contenido (CMS), perfiles de usuario.
  \item Eventos/logs y telemetría (con diseño adecuado e índices).
  \item APIs con payload JSON y modelo agregado (documentos anidados).
  \item Microservicios con datos por dominio (bounded context).
\end{itemize}
\end{frame}

\begin{frame}[fragile]{MongoDB: programas y comandos clave}
\textbf{Programas/herramientas:}
\begin{itemize}
  \item \texttt{mongod} (servidor)
  \item \texttt{mongosh} (shell/CLI)
  \item \texttt{mongodump}/\texttt{mongorestore} (backup/restore)
  \item \texttt{mongoexport}/\texttt{mongoimport} (CSV/JSON)
  \item MongoDB Compass (GUI)
\end{itemize}

\textbf{Comandos típicos:}
\begin{verbatim}
mongod --dbpath /data/db
mongosh "mongodb://localhost:27017"

mongodump --db tienda --out backup/
mongorestore --db tienda backup/tienda
\end{verbatim}
\end{frame}

\begin{frame}[fragile]{MongoDB: operaciones esenciales (ejemplos)}
\begin{verbatim}
// Seleccionar BD y colección
use tienda
db.createCollection("productos")

// Insertar
db.productos.insertOne({ sku:"A1", nombre:"Teclado",
  precio:29.9, tags:["pc","periferico"] })

// Consultar con filtro y proyección
db.productos.find(
  { precio: { $lt: 50 } },
  { nombre:1, precio:1, _id:0 }
)

// Indices
db.productos.createIndex({ sku: 1 }, { unique: true })
db.productos.createIndex({ tags: 1 })
\end{verbatim}
\end{frame}

\begin{frame}[fragile]{MongoDB: Aggregation Pipeline (ejemplo)}
\begin{verbatim}
db.ventas.aggregate([
  { $match: { fecha: { $gte: ISODate("2026-01-01") } } },
  { $group: { _id: "$sku", total: { $sum: "$importe" }, n: { $sum: 1 } } },
  { $sort: { total: -1 } },
  { $limit: 10 }
])
\end{verbatim}
\end{frame}

% =========================================================
\section{Cierre}

\begin{frame}{Comparativa rápida: ¿cuándo usar cuál?}
\begin{itemize}
  \item \textbf{SQLite}: simple, local/embebido, transacciones, una app/un dispositivo.
  \item \textbf{Neo4j}: relaciones complejas, caminos, recomendación, fraude, knowledge graph.
  \item \textbf{MongoDB}: documentos JSON, esquema flexible, agregaciones, APIs y escalado.
\end{itemize}
\end{frame}

\begin{frame}{Checklist de decisión}
\begin{itemize}
  \item ¿Tus consultas principales son \textbf{joins} relacionales clásicos? $\rightarrow$ SQLite (o SQL en general).
  \item ¿Tu valor está en \textbf{relaciones y caminos}? $\rightarrow$ Neo4j.
  \item ¿Tu modelo natural es \textbf{documentos JSON} con evolución de campos? $\rightarrow$ MongoDB.
  \item ¿Necesitas simplicidad operativa local? $\rightarrow$ SQLite.
\end{itemize}
\end{frame}

\end{document}
